% Figure IV.10 --- assurance position, not a prose matrix; position is the
% one channel a reader reads accurately (Cleveland-McGill).
%
% Reader question: which threats are mechanized, which are only bounded
%   under a stated model, and which remain open?
% Claim: assurance advances only when a threat acquires a concrete artifact
%   or theorem; most rows are conditional on a named hypothesis, and three
%   remain open.
% Evidence: the seven threat rows and their own evidence phrase (signed
%   card, the Theta(log m) bound, the settlement gate, ...) from S:fh-limitations.
% Must distinguish: OPEN (acknowledged gap) from CONDITIONAL (bound holds
%   under a named hypothesis) from MECHANIZED (artifact exists).
% Grammar chosen: position on one shared axis, seven rows (a banded ledger).
%   Grammar rejected: a checklist table loses the ORDERING claim (rows are
%   not equally mechanized); a decorative shield/lock icon per row asserts
%   security the paper does not claim.
% Five-second test: which single row is closest to MECHANIZED, and which
%   three rows have not moved off OPEN?
\pdfigurehue{pdgold}%
\begin{figure*}[t]
\centering
\begin{tikzpicture}[pd figure,x=1cm,y=1cm]
  \def\lx{1.65}   \def\x0{2.05}   \def\pw{8.35}
  \node[pd title,anchor=north west,text width=10.6cm,align=left] at (0,5.15)
    {assurance advances only when a threat acquires an artifact or theorem};

  \foreach \y/\lab in {3.30/{identity + replay},1.90/{revocation},0.50/{equivocation},
    -0.90/{settlement},-2.30/{cold-start admission},-3.70/{cartel monitoring},
    -5.40/{centralization}}
    {\node[pd label,anchor=east,text width=2.7cm,align=right] at (\lx,\y) {\lab};
     \draw[pd hairline] (\x0,\y) -- (\x0+\pw,\y);}

  % identity + replay -- strongest evidence: chapter's own exemplar
  \draw[pd focus rule,line width=1.2pt] (\x0,3.30) -- (10.05,3.30);
  \node[pd focus datum] at (10.05,3.30) {};
  \node[pd focus label,anchor=south east,text width=3.4cm,align=right] at (10.05,3.48)
    {signed card $+$ epoch root};

  % revocation -- literally the concept; breach hue
  \draw[pd breach rule,line width=1.2pt] (\x0,1.90) -- (7.27,1.90);
  \node[pd breach datum] at (7.27,1.90) {};
  \node[pd breach label,anchor=south east,text width=3.6cm,align=right] at (7.27,2.08)
    {$\Theta(\log m)$ only under reliable rounds};

  % equivocation -- conditional, neutral ink
  \draw[pd rule,line width=1.2pt] (\x0,0.50) -- (6.92,0.50);
  \node[pd datum] at (6.92,0.50) {};
  \node[pd label,anchor=south east,text width=3.4cm,align=right] at (6.92,0.68)
    {proof after roots co-locate};

  % settlement -- wide uncertainty interval, neutral ink (thick = range)
  \draw[pd rule,line width=1.2pt] (\x0,-0.90) -- (6.75,-0.90);
  \draw[pd rule,line width=2.2pt] (5.09,-0.90) -- (7.35,-0.90);
  \node[pd datum] at (6.75,-0.90) {};
  \node[pd label,anchor=south east,text width=3.6cm,align=right] at (6.75,-0.72)
    {bypassable to gated: gate-dependent};

  % cold-start admission -- open, neutral ink
  \draw[pd rule,line width=1.2pt] (\x0,-2.30) -- (4.96,-2.30);
  \node[pd datum] at (4.96,-2.30) {};
  \node[pd label,anchor=south west,text width=4.0cm,align=left] at (2.45,-2.12)
    {thin reputation cannot price a bond};

  % cartel monitoring -- open, neutral ink
  \draw[pd rule,line width=1.2pt] (\x0,-3.70) -- (4.66,-3.70);
  \node[pd datum] at (4.66,-3.70) {};
  \node[pd label,anchor=south west,text width=4.0cm,align=left] at (2.45,-3.52)
    {monitoring game incomplete};

  % centralization -- open, neutral ink; the pushback point (S:fh-limitations)
  \draw[pd rule,line width=1.2pt] (\x0,-5.40) -- (4.22,-5.40);
  \node[pd datum] at (4.22,-5.40) {};
  \node[pd label,anchor=south west,text width=4.0cm,align=left] at (2.45,-5.22)
    {pushback is design discipline, not a bound};

  % ---- assurance-position axis --------------------------------------------
  \draw[pd rule,-{Stealth[length=1.8mm,width=1.1mm]}] (\x0,-6.40) -- (\x0+\pw+0.05,-6.40);
  \foreach \x/\lab in {2.05/OPEN,6.15/CONDITIONAL,10.25/MECHANIZED}
    {\draw[pd tick] (\x,-6.40) -- (\x,-6.55);
     \node[pd label,anchor=north] at (\x,-6.60) {\lab};}
  \node[pd note,anchor=north west,text width=10.2cm] at (0,-7.05)
    {position, not colour, carries the claim: each row advances only when the
    missing artifact or theorem exists};
\end{tikzpicture}
\caption{\textbf{The assurance landscape is uneven and falsifiable.} Position
encodes how far each threat has moved from an acknowledged gap toward a
mechanized artifact. Identity and replay (gold) have the strongest concrete
evidence. Revocation (red) is conditional on connected reliable rounds
(Theorem~\ref{thm:fh-conv}); equivocation and settlement are conditional on
named connectivity or custody hypotheses, and the wide settlement interval
shows how far its assurance moves if the custody gate is bypassable
(\S\ref{sec:fh-escrow}). Cold-start admission, cartel monitoring, and
centralization remain open. A row advances only when the missing artifact or
theorem exists. [internal]}
\label{fig:fh-threat-bands}
\end{figure*}
